\scanpage{33}
Sous \uline{ces conditions}, \uline{$F$ est représentable}.
\margin{Transformer $4^\circ)$ sous une forme plus facile à vérifier, avec un bon \uline{\unsure{énoncé-type} de dim 1}}

En effet, on peut recoller $X_1$ et $X'$ suivant $u':X'_1\mathbin{-}X'$, on obtient $X$, et sur $X$ grâce à $\underset{\leftarrow}{L}_1$ on a un $\xi\in F(X)$, induisant $\xi_1$ et $\xi'_1$. On lui applique le corollaire 1, qui indique que \deleted{il} \inserted{$(X,\xi)$} représente $F$.

\figure{\begin{tikzpicture}[x=1cm,y=1cm]
\draw (0,0)--(.85,0);\draw (.17,-.1)--(.68,-.1);\draw (.27,-.19)--(.56,-.19);
\end{tikzpicture}}

\uline{Prop. 3} Soit $F:(\mathrm{Sch})_{/S}^{\circ}\mathbin{-}\mathrm{Ens}$, on suppose: \deleted{que}

$1^\circ)$ Pour $\forall a\in S$, posant $S_a=\operatorname{Spec}\mathcal O_{S,a}$, le foncteur $F_a=F_{/S_a}$ est représentable par un préschéma \inserted{$X_a$} loc. de t. f. sur $S_a$.

$2^\circ)$ Pour $\forall T$ sur $S$ loc. de type fini sur $S$, et $\eta\in F(T)$, l'ensemble \inserted{$M(T,\eta)$} $t\in T$ tels que $\eta_t\in F(\unsure{T_a|T_t})$ soit \inserted{\uline{loc}} \uline{modulaire}, est une partie \uline{constructible} de $T$.
\margin{\uline{Transformer en une condition de modularité ``générique''}}
\inserted{$(3^\circ)$ $F$ satisfait $\underset{\leftarrow}{L}_1$ et $\underset{\leftarrow}{L}_2$}

Sous ces conditions, $F$ est représentable.
\mspage{32}
